\documentclass[uplatex,a4paper,11pt]{jsarticle}

% ===== パッケージ =====
\usepackage[dvipdfmx]{graphicx}
\usepackage{amsmath, amssymb, amsthm}
\usepackage{geometry}
\usepackage{tikz}
\usetikzlibrary{automata, positioning}
\usepackage{hyperref}
\usepackage{here}
\usepackage{bookmark}

\geometry{margin=25mm}
\renewcommand{\baselinestretch}{1.3}

% ===== 定理環境 =====
\newtheorem{theorem}{定理}

% ===== コマンド =====
\newcommand{\N}{\mathbb{N}}

% ===== タイトル =====
\title{Homework 1b}
\author{学籍番号: 1270374 \\ 名前: 松本 奈生}
\date{\today}

\begin{document}

\maketitle

% =========================
% Homework 1b
% =========================

\section*{Homework 1b}

\section*{Question 1 (Perfect Matching と CNF)}

\subsection*{証明}

グラフ $G=(V,E)$ に対して，$\phi_G$ が充足可能であることと，
$G$ が完全マッチングを持つことが同値であることを示す．

\subsubsection*{(i) $G$ が完全マッチングを持つ $\Rightarrow$ $\phi_G$ は充足可能}

完全マッチング $M$ を用いて，

\[
\theta(x_e)=
\begin{cases}
1 & (e \in M)\\
0 & (e \notin M)
\end{cases}
\]

と定める．

各頂点にはちょうど1本の辺が接続されるため，
$\phi_{1,G}$ と $\phi_{2,G}$ の両方を満たす．

よって $\phi_G$ は充足可能．

\subsubsection*{(ii) $\phi_G$ が充足可能 $\Rightarrow$ 完全マッチングが存在}

\[
M = \{e \in E \mid \theta(x_e)=1\}
\]

とする．

各頂点について「少なくとも1本」かつ「高々1本」なので，
ちょうど1本選ばれる．

よって $M$ は完全マッチング．

\[
\boxed{\text{証明完了}}
\]

% =========================

\section*{Question 2 (4正則グラフ)}

\subsection*{証明}

$n \ge 5$ とする．  
頂点 $V=\{0,1,\dots,n-1\}$ を円周上に配置する．

\[
E = \{\{i,i+1\}, \{i,i+2\} \mid i \in V\} \ (\mathrm{mod}\ n)
\]

各頂点は $i\pm1, i\pm2$ に接続するので，

\[
\deg(i)=4
\]

よって4正則グラフが存在する．


% =========================

\section*{Question 3 (Petersenグラフの彩色数)}

\subsection*{証明}

Petersenグラフの彩色数 $\chi(G)$ を求めるため，
必要な色数について場合分けして考える．

まず，$\chi(G) < 3$ の場合を考える．  
これは $\chi(G) \le 2$，すなわち2色で彩色できることを意味する．  
しかし，2色で彩色できるグラフは二部グラフである．

一方，Petersenグラフは長さ5のサイクル（奇数長サイクル）を含むため，
二部グラフではない．  
したがって，2色での彩色は不可能であり，

\[
\chi(G) \ge 3
\]

が成り立つ．

次に，$\chi(G) = 3$ の場合を考える．  
実際にPetersenグラフは3色で彩色することができる．
例えば，外側の5頂点を2色で交互に塗り，
内側の頂点を適切に3色目を含めて塗ることで，
隣接する頂点が異なる色になるようにできる．

したがって，

\[
\chi(G) \le 3
\]

が成り立つ．

以上より，

\[
\chi(G) = 3
\]

である．

% =========================

\section*{Question 4}
\subsection*{解答}

PetersenグラフがHamiltonパスを持つが，Hamilton閉路を持たないことを示す．

まず，Hamiltonパスの存在について考える．
Petersenグラフは対称性の高いグラフであるため，適切に頂点を選べば，
すべての頂点を一度ずつ通る経路が存在すると考えられる．

実際に頂点を順に辿っていくと，隣接する頂点をうまく選ぶことで，
全頂点を重複なく通ることができる．
したがって，PetersenグラフはHamiltonパスを持つ．

次に，Hamilton閉路の存在について考える．ここでは背理法を用いる．
PetersenグラフがHamilton閉路を持つと仮定する．

Petersenグラフは3正則グラフであり，各頂点の次数は3である．
Hamilton閉路では，各頂点はちょうど2本の辺で接続されるため，
各頂点には1本の使われない辺が残る．

この残った辺を集めると，各頂点にちょうど1本ずつ対応するので，
それらは完全マッチングを構成することになる．

ここで，Petersenグラフの構造を考えると，
このような完全マッチングを取り除いた場合，
残りのグラフは1つの閉路とはならず，
複数の連結成分に分かれてしまう．

しかし，Hamilton閉路が存在するならば，
残りの辺は1つの閉路を構成するはずであり，これは矛盾である．

したがって，PetersenグラフはHamilton閉路を持たない．

以上より，PetersenグラフはHamiltonパスを持つが，
Hamilton閉路は持たないことが示された．

% =========================

\section*{Question 5}

\subsection*{証明}

頂点数 $n \ge 2$ のグラフを考える．

次数は $0$ から $n-1$ の範囲であるが，

- 次数0があると次数$n-1$は存在できない

したがって次数の種類は高々 $n-1$ 個．

頂点は $n$ 個あるので，鳩の巣原理より  
同じ次数を持つ頂点が存在する．


\end{document}

